\documentclass[11pt,a4paper]{article}


% ══════════════════════════════════════════════════════════════════════════════
%
% Packages
%

\usepackage[margin=2.5cm]{geometry}
\usepackage[utf8]{inputenc}
\usepackage[T2A,T1]{fontenc}
\usepackage[bulgarian]{babel}
\usepackage{lmodern}
\usepackage{microtype}
\usepackage{amsmath,amssymb}
\usepackage{graphicx}
\usepackage{booktabs}
\usepackage{multirow}
\usepackage{array}
\usepackage{hyperref}
\usepackage{xcolor}
\usepackage{titlesec}
\usepackage{abstract}
\usepackage{parskip}
\usepackage{enumitem}
\usepackage{algorithm}
\usepackage{algpseudocode}
\usepackage{caption}
\usepackage{subcaption}
\usepackage{natbib}
\usepackage{xspace}

\hypersetup{
    colorlinks=true,
    linkcolor=blue!60!black,
    citecolor=green!50!black,
    urlcolor=blue!70!black
}

\titlespacing*{\section}{0pt}{1.4ex plus 0.5ex minus 0.2ex}{0.8ex plus 0.2ex}
\titlespacing*{\subsection}{0pt}{1.2ex plus 0.4ex minus 0.2ex}{0.6ex plus 0.2ex}

\newcommand{\todo}[1]{\textcolor{red}{[TODO: #1]}}
\newcommand{\eg}{{напр.}\xspace}
\newcommand{\ie}{{т.е.}\xspace}

\begin{document}


% ══════════════════════════════════════════════════════════════════════════════
%
% Заглавие
%

\title{Дълбоко обучение с подкрепление за настолна игра "Заселнците на Катан"}

\author{
    Александър Македонски \\
    Факултет по математика и информатика \\
    Софийски университет „Св. Климент Охридски" \\
    {alemak2002@gmail.com} \\
}
\date{24.06.2026}
\maketitle

% ══════════════════════════════════════════════════════════════════════════════
%
% Тяло
%
%

\begin{abstract} % 
Представям система за дълбоко обучение с подкрепление (DRL), предназначена
да се научи да играе Заселниците на Катан --- сложна многоиграчова
настолна игра, която комбинира стохастично генериране на ресурси, преговори и
дългосрочно стратегическо планиране. Имплементирам и сравнявам три алгоритъма
за DRL --- Deep Q-Network (DQN), Proximal Policy Optimisation (PPO) с
Generalised Advantage Estimation (GAE) и базова линия REINFORCE --- използвайки
средата catanatron gym. Ключов принос е плътна, оформена функция за
награда, която декомпозира оскъдния краен сигнал на победа в
стъпкови междинни награди за построяване на пътища, селища и градове, както и
термин за поставяне на разбойника, който награждава враждебността към опоненти.
Моят най-добър агент --- Dueling Double DQN (D3QN), трениран чрез
самостоятелна игра --- постига 90\% на победи срещу базовата
линия на (претеглен) случаен играч. Стандартен DQN с оформена награда достига
85\%, а PPO с бинарна награда победа/загуба --- 76\%.
Aнализираме влиянието им върху тренировката. Допълнително: Уеб-базиран UI
позволява на хора да се изправят срещу тренираните модели. Отворен код. 
\end{abstract}

% ══════════════════════════════════════════════════════════════════════════════
\section{Въведение}

Настолните игри отдавна служат за ориентир в изследванията в
областта на изкуствения интелект --- от дама и
шах~\citep{samuel1959some,silver2017mastering} до
Го~\citep{silver2016mastering}. Заселниците на Катан представлява
особено интересно предизвикателство, тъй като съчетава {стохастични
елементи} (хвърляне на зарове на всеки ход, случайна начална наредба на средата),
{голямо дискретно пространство от действия} (над 300 допустими действия
на стъпка), {дълги епизоди} (стотици ходове на игра) и
{многоиграчова конкурентна} динамика. Тези свойства правят стандартни DRL подходи нетривиални.

Съществуващите изследвания за ИИ в Катан са разглеждали Monte Carlo Tree
Search и евристики, базирани на
правила~\citep{szita2010monte}, но подходите с дълбоко RL остават
слабо проучени.

В настоящата работа изследваме дали съвременните DRL алгоритми могат да
научат конкурентна игра на Катан от нулата. Нашите основни приноси са:

\begin{enumerate}[leftmargin=1.5em]
    \item Плътна функция за оформяне на наградата, която декомпозира сигнала
    за победа в интерпретируеми компоненти за отделни събития, включително нов
    термин за поставяне на разбойника, който различава нараняването на опоненти
    от наказването на самия агент.
    \item Емпирично сравнение на DQN (стандартен и Dueling), PPO с GAE и
    REINFORCE в двуиграчова постановка на Катан.
    \item Анализ на това как тренировката чрез самостоятелна игра (сам срещу себе си) подобрява
    представянето на DQN.
    \item Дпълнение - Уеб-базиран UI за игра човек срещу агент, изграден върху
    визуализатора на {catanatron}.
\end{enumerate}

% ══════════════════════════════════════════════════════════════════════════════
\section{Свързани изследвания}

\paragraph{ИИ за Катан.}
\citet{szita2010monte} прилагат Monte Carlo Tree Search за Катан, постигайки
приемлива игра, но разчитайки на ръчно изградена политика за разгръщане.
\citet{pfeiffer2004} изследват агенти, базирани на правила. По-наскоро
рамката {catanatron} \citep{collazo2021catanatron} предоставя среда,
съвместима с gym, и базова линия с претеглен случаен играч. По наши данни
никое предходно изследване не е публикувало резултати от DRL тренировка от
край до край в многоиграчовата gym постановка.

\paragraph{Deep Q-Networks.}
Алгоритъмът DQN~\citep{mnih2015human} - повторение на опита, целева
мрежа за стабилизиране. Dueling DQN~\citep{wang2016dueling} --- стойност на
състоянието $V(s)$ и предимство с нулево средно $A(s,a)$, което е особено
полезно когато много действия имат подобна очаквана възвращаемост. Double DQN~\citep{hasselt2016deep} - адресира надценяването на
Q-стойностите.

\paragraph{Методи на градиента на политиката.}
PPO~\citep{schulman2017proximal} --- предлага и по-стабилна тренировка от
обикновения градиент на политиката. GAE~\citep{schulman2015high} интерполира между
Монте Карло оценки и едностъпкови TD оценки, подходящ за дълги игри.

\paragraph{Самостоятелна игра.}
Ние приемаме опростена версия, при
която замразено копие на текущия най-добър модел служи като опонент, позволявайки
на агента да се подобрява срещу нестационарен, но все по-силен противник.

% ══════════════════════════════════════════════════════════════════════════════
\section{Метод}

\subsection{Среда}

Използваме gym средата {catanatron} ({catanatron/Catanatron-v0}).
Наблюдението е плосък вектор с размер $194N + 226$, където $N$ е броят
играчи. За двуиграчова игра това дава 614-мерно наблюдение. Пространството от
действия е дискретно множество от $|\mathcal{A}| = 328$ цели числа. На всяка
стъпка средата връща и множество от {допустими действия}; недопустимите
действия се маскират с малка стойност, за да не може агентът да ги избере.

\subsection{Оформена функция за награда}

Наградата по подразбиране на средата е бинарна (+1 за победа, $-$1 за загуба).
Дефинираме междинна оформена награда $r_t$, която допълва този сигнал:

\begin{equation}
r_t = r_t^{\text{точки}} + r_t^{\text{път}} + r_t^{\text{град}} +
      r_t^{\text{селище}} + r_t^{\text{победа}}
\end{equation}

\noindent където всеки термин е дефиниран по следния начин. ($\Delta$ означава разликата между сегашното наблюдение и предходното.)\
$
r_t^{\text{точки}} = \max\bigl(0,\;\Delta\text{точки}_t\bigr) \cdot c_{\text{VP}},\ \text{където}\  c_{\text{точки}} = 3\ 
$
Награждават се само положителни печалби на ТП --- временни загуби ($\eg$ загуба на Най-дълъг път) не се наказват, за да не се обезкуражи строенето на пътища;
$
r_t^{\text{път}} = \max(0, \Delta\text{пътища}_t)   \cdot c_{\text{път}},\ \text{където}\  c_{\text{път}} = 0,2;\ 
r_t^{\text{град}} = \max(0, \Delta\text{градове}_t)  \cdot c_{\text{град}},\ \text{където}\  c_{\text{град}} = 1,5;\  
r_t^{\text{селище}} = \max(0, \Delta\text{селища}_t - 2) \cdot c_{\text{селище}},\ \text{където}\  c_{\text{селище}} = 1;\  
$ 
Награждават се само селища, построени след задължителните две начални; и

$
r_t^{\text{победа}} = \begin{cases} +20 & \text{ако агентът печели на стъпка } t \\ -20 & \text{ако опонентът печели на стъпка } t \\ 0 & \text{иначе} \end{cases}
$

По-късно се представя и друга награда, която за съжаление не влиза в тестовете --- $r_t^{\text{разбойник}}$:

\paragraph{Поставяне на разбойника.}
Когато агентът поставя разбойника, оценяваме качеството на избраната плочка:
\begin{equation}
r_t^{\text{разбойник}} = c_{\text{р}} \sum_{v \in \mathcal{N}}
\begin{cases}
  c_{\text{опонент}} \cdot f_v & \text{ако сградата при } v \text{ е на опонент} \\
  c_{\text{собствено}}    \cdot f_v & \text{ако сградата при } v \text{ е на агента}
\end{cases}
\end{equation}
където $\mathcal{N}$ е множеството от шестте върха около
плочката с разбойника, $f_v = c_{\text{град\_фактор}}$ ако има град и 1 в
противен случай, $c_{\text{р}} = 0{,}4$, $c_{\text{опонент}} = 1{,}2$,
$c_{\text{собствено}} = -1$ и $c_{\text{град\_фактор}} = 2$. Асиметрията
$c_{\text{опонент}} > c_{\text{собствено}}$ предава идеята, че нараняването на
опонент е по-ценно от защитата на себе си.

\subsection{Архитектура на мрежите}
Всички мрежи използват една и съща начална аргитектура:
\[
\mathbf{h} = \text{ReLU}\bigl(W_3\,\text{ReLU}(W_2\,\text{ReLU}(W_1 \mathbf{s}))\bigr)
\]
Някои модели имплементират доълнителни слоеве над нея или я разширяват с по-голям брой неврони (в повечето случаи двоен).

\subsection{Архитектура на Dueling DQN}
Вариантът Dueling DQN разделя крайното представяне на два клона:
\begin{equation}
Q(s, a) = V(s) + \Bigl(A(s, a) - \frac{1}{|\mathcal{A}|}\sum_{a'}A(s,a')\Bigr)
\end{equation}
Тoва помага, когато много действия имат подобно предимство, което
е честа ситуация в Катан (\eg 50\% от състоянията са „край на хода", които дават 0 точки).

За тренировка се изпoлзва загубата на Хубер между предсказаните и целевите
Q-стойности:
\begin{equation}
\mathcal{L} = \mathbb{E}_{(s,a,r,s') \sim \mathcal{D}}
\Bigl[\delta\!\left(Q(s,a) - \Bigl(r + \gamma \max_{a'}\hat{Q}(s',a')\Bigr)\right)\Bigr]
\end{equation}
където $\hat{Q}$ е целевата мрежа и $\delta(\cdot)$ е загубата на Хубер.
Недопустимите действия се маскират до $-\infty$ преди операцията $\max$ за избиране на действие в средата.

\subsection{PPO с GAE}

Актьорът и критикът на PPO са отделни MLP-та, споделящи същата скрита
архитектура. Отрязаната сурогатна цел е:
\begin{equation}
\mathcal{L}^{\text{CLIP}} = \mathbb{E}_t\Bigl[
  \min\bigl(r_t(\theta)\hat{A}_t,\;
  \text{clip}(r_t(\theta), 1{-}\epsilon, 1{+}\epsilon)\hat{A}_t\bigr)
\Bigr]
\end{equation}
Предимствата $\hat{A}_t$ се оценяват чрез GAE с $\lambda \in \{0{,}90, 0{,}95\}$
(претърсвано по време на тренировка).

\subsection{Самостоятелна игра}

Имплементираме самостоятелна игра за DQN, където агентът е изправен срещу свое копие, което се обновява периодично. Така агентът се подобрява заедно с опонента, предотвратявайки прекалено нагаждане към фиксираната базова линия.

% ══════════════════════════════════════════════════════════════════════════════
\section{Експериментална постановка}

\subsection{Среда и базови линии}

Всички експерименти използват {catanatron/Catanatron-v0} gym в
двуиграчова конфигурация. Основната базова линия е WeightedRandomPlayer
от {catanatron}, който избира действия с ръчно настроени евристични
тегла. Равномерен RandomPlayer се използва като по-слаба референция.

\subsection{Данни и събиране на епизоди}

Няма офлайн набор от данни; всички тренировъчни данни се генерират онлайн чрез
самостоятелна игра или игра срещу базови линии. Всеки епизод съответства на
една пълна игра на Катан (обикновено 100--400 хода). Тренировъчите цикли
варират от 5\,000 до 20\,000 епизода.

\subsection{Хиперпараметри}

\begin{table}[h]
\centering
\caption{Ключови хиперпараметри за всеки алгоритъм.}
\label{tab:hparams}
\begin{tabular}{lll}
\toprule
Параметър & DQN / D3QN & PPO \\
\midrule
Оптимизатор             & Adam              & Adam \\
Скорост на учене        & $1 \times 10^{-4}$ & $3 \times 10^{-4}$ \\
Отстъпка $\gamma$       & 0{,}99            & 0{,}99 \\
Размер на пакет         & 128               & пълен rollout \\
Памет за повторение     & 100\,000          & --- \\
$\epsilon$ начало/край  & 1{,}0 / 0{,}05    & --- \\
Затихване на $\epsilon$ & $1 \times 10^{-6}$/стъпка & --- \\
Синхр. на целева мрежа  & на всеки 5 епизода & --- \\
GAE $\lambda$           & ---               & 0{,}95 \\
PPO clip $\epsilon$     & ---               & 0{,}2 \\
Коеф. на ентропия       & ---               & 0{,}01 \\
\bottomrule
\end{tabular}
\end{table}


\subsection{Метрика за оценка}

Основната метрика е процентът на победи за 500 оценъчни епизода
срещу WeightedRandomPlayer, изиграни след тренировката. Вторичните метрики са
средната дължина на играта (брой ходове), общата награда за епизод и средната
загуба за епизод.

Създадени са и доълнителни оценки за игри изиграни срещу фиксирани играчи от catanatron (Alpha Beta Pruning или MCTS алгоритми.)

\subsection{Хардуер и софтуер}

Експериментите са проведени на единична потребителска графична карта. Всички модели са
имплементирани в PyTorch. Средата е catanatron v$\geq$1.0
\citep{collazo2021catanatron}, а взаимодействията с gym използват API-я на
{gymnasium}. Тренировъчните статистики се записват като JSON
и се визуализират чрез персонализирано Streamlit табло.

% ══════════════════════════════════════════════════════════════════════════════
\section{Резултати}

\subsection{Сравнение на процента на победи}
\begin{table}[H]
\centering
\caption{Процент на победи срещу WeightedRandomPlayer (двуиграчова игра).
         По-висок е по-добър; 50\% съответства на игра на ниво случаен избор.}
\label{tab:results}
\begin{tabular}{llcc}
\toprule
Алгоритъм & Награда & Епизоди & Победи (\%) \\
\midrule
Случайна базова линия    & ---               & ---     & $\sim$50 \\
DQN (стандартен)         & Оформена          & 10\,000 & 85 \\
D3QN + самост. игра      & Оформена          & 15\,000 & 90 \\
PPO + GAE                & Победа/загуба ($\pm$1) & 20\,000 & 76 \\
REINFORCE                & Оформена          & 10\,000 & 25 \\
\bottomrule
\end{tabular}
\end{table}

\subsection{Ефект на самостоятелната игра}

Замяната на фиксирания опонент WeightedRandomPlayer с замразено копие на
текущия DQN модел увеличава крайния процент на победи от 85\% до 90\% при
15\,000 епизода. Това предполага, че адаптивната програма, генерирана от
самостоятелната игра, осигурява по-трудни и разнообразни опоненти, подобрявайки
обобщаването.

% ══════════════════════════════════════════════════════════════════════════════
\section{Аблационни изследвания}

\begin{table}[h]
\centering
\caption{Аблационно изследване на компонентите на наградата за DQN (10\,000 епизода).}
\label{tab:ablation}
\begin{tabular}{lc}
\toprule
Конфигурация & Победи (\%) \\
\midrule
Пълна оформена награда                           & 85 \\
\quad $-$ награда за пътища                      & +2 пп () \\
\quad $-$ термини за градове/селища              & +10 пп () \\
\quad $-$ термин за разбойника                   & все още не е тествано \\
Само бинарна победа/загуба                       & $\approx$50 (без подобрение) \\
Стандартен DQN $\to$ Dueling DQN                 & +5 (от 85 до 90 при самост. игра) \\
Фиксиран опонент $\to$ самостоятелна игра        & +5 (от 85 до 90) \\
\bottomrule
\end{tabular}
\end{table}

Преминаването от стандартен DQN към Dueling архитектура подобри стабилността
на тренировката, в съответствие с хипотезата, че много действия в Катан имат
подобни стойности на награда и оценката на състоянието трябва
да бъде отделена. Бинарната награда победа/загуба самостоятелно беше
недостатъчна: тя не доведе до подобрение, потвърждавайки необходимостта от плътно оформяне при дълги игри.

% ══════════════════════════════════════════════════════════════════════════════
\section{Дискусия и ограничения}
\label{sec:bugs}

\subsection{Критични бъгове}

Поради непредвидени бъгове в ръчно написан код, се отдели много време в трениране на безполезни агенти. Дори при писане на тестове на кода, бъгът остана скрит, тъй като беше недетерминистичен. За повече информация, може да се види дневника на проекта, където са описани всички стъпки на развитиети на проекта.

\subsection{Ограничения}

\begin{itemize}[leftmargin=1.5em]
    \item Всички агенти се оценяват в двуиграчова постановка.
    Четирииграчовата постановка беше изследвана накратко.
    \item Векторът на наблюдение е плосък масив от характеритики,
    който не кодира изрично пространствената структура на дъската. Граф
    невронна мрежа или CNN върху тензора на дъската биха уловили локалността
    по-добре, но този метод не бе изследван.
    \item REINFORCE се представи изключително слабо; не изследвахме
    напълно дали това се дължи на бъг в тренировката или на алгоритмично
    ограничение.
    \item Търсенето на хиперпараметри беше ръчно и ограничено. Систематично
    претърсване (\eg Байесова оптимизация) би могло да даде допълнителни
    подобрения.
    \item Скоростта на тренировка беше: приблизително 2\,000
    епизода за 30 минути на единична GPU.
\end{itemize}

% ══════════════════════════════════════════════════════════════════════════════
\section{Бъдеща работа.}
Обещаващи посоки включват (1) прилагане на граф невронни мрежи върху
графовата структура на дъската, (2) пълна четирииграчова самостоятелна игра
с population-based тренировка, (3) интегриране на вероятността за производство
на ресурси в наблюдението за подобряване на решенията за поставяне на
селища, и (4) система за търговия между играчите, изпълняваща се от отделна невронна мрежа.
\\
\\

% ══════════════════════════════════════════════════════════════════════════════
\bibliographystyle{plainnat}
\begin{thebibliography}{99}

\bibitem[Collazo(2021)]{collazo2021catanatron}
Collazo, B. (2021).
\newblock Catanatron: A Catan engine and OpenAI gym environment.
\newblock \url{https://github.com/bcollazo/catanatron}.

\bibitem[Hasselt et al.(2016)]{hasselt2016deep}
van Hasselt, H., Guez, A., and Silver, D. (2016).
\newblock Deep reinforcement learning with double Q-learning.
\newblock In {AAAI}, pages 2094--2100.

\bibitem[Mnih et al.(2015)]{mnih2015human}
Mnih, V., Kavukcuoglu, K., Silver, D., et al. (2015).
\newblock Human-level control through deep reinforcement learning.
\newblock {Nature}, 518(7540):529--533.

\bibitem[Pfeiffer(2004)]{pfeiffer2004}
Pfeiffer, M. (2004).
\newblock Reinforcement learning of strategies for Settlers of Catan.
\newblock In {Proceedings of the International Conference on Computer
Games}.

\bibitem[Samuel(1959)]{samuel1959some}
Samuel, A.~L. (1959).
\newblock Some studies in machine learning using the game of checkers.
\newblock {IBM Journal of Research and Development}, 3(3):210--229.

\bibitem[Schulman et al.(2015)]{schulman2015high}
Schulman, J., Moritz, P., Levine, S., Jordan, M., and Abbeel, P. (2015).
\newblock High-dimensional continuous control using generalized advantage
estimation.
\newblock {arXiv preprint arXiv:1506.02438}.

\bibitem[Schulman et al.(2017)]{schulman2017proximal}
Schulman, J., Wolski, F., Dhariwal, P., Radford, A., and Klimov, O. (2017).
\newblock Proximal policy optimization algorithms.
\newblock {arXiv preprint arXiv:1707.06347}.

\bibitem[Silver et al.(2016)]{silver2016mastering}
Silver, D., Huang, A., Maddison, C.~J., et al. (2016).
\newblock Mastering the game of Go with deep neural networks and tree search.
\newblock {Nature}, 529(7587):484--489.

\bibitem[Silver et al.(2017)]{silver2017mastering}
Silver, D., Schrittwieser, J., Simonyan, K., et al. (2017).
\newblock Mastering chess and shogi by self-play with a general reinforcement
learning algorithm.
\newblock {arXiv preprint arXiv:1712.01815}.

\bibitem[Szita et al.(2010)]{szita2010monte}
Szita, I., Chaslot, G., and Spronck, P. (2010).
\newblock Monte-Carlo tree search in Settlers of Catan.
\newblock In {Advances in Computer Games}, pages 21--32.

\bibitem[Wang et al.(2016)]{wang2016dueling}
Wang, Z., Schaul, T., Hessel, M., Lanctot, M., Freitas, N., and Silver, D.
(2016).
\newblock Dueling network architectures for deep reinforcement learning.
\newblock In {ICML}, pages 1995--2003.

\end{thebibliography}

\end{document}
